% =====================================================================
% DBA - Proyecto final - Entrega 1
% Diseño del esquema NoSQL y plan de implementación
% =====================================================================
\documentclass[11pt,a4paper]{article}
% =====================================================================
% Preámbulo compartido - DBA Proyecto final
% =====================================================================
\usepackage[utf8]{inputenc}
\usepackage[T1]{fontenc}
\IfFileExists{lmodern.sty}{\usepackage{lmodern}}{}
\usepackage[spanish,es-noshorthands,es-tabla]{babel}
\usepackage[margin=2.5cm]{geometry}
\usepackage{amsmath,amssymb}
\usepackage{booktabs,tabularx,array}
\usepackage{enumitem}
\usepackage[dvipsnames,table]{xcolor}
\usepackage{tikz}
\usetikzlibrary{calc,positioning}
\usepackage{listings}
\usepackage{caption}
\usepackage{float}
\usepackage[hidelinks]{hyperref}
\usepackage{xurl}

% --- Colores --------------------------------------------------------
\definecolor{acento}{HTML}{2F6FDB}
\definecolor{gris}{HTML}{6B6B6B}
\hypersetup{colorlinks=true, linkcolor=acento, urlcolor=acento, citecolor=acento}

% --- Tablas ---------------------------------------------------------
\newcolumntype{L}{>{\raggedright\arraybackslash}X}
\renewcommand{\arraystretch}{1.25}

% --- Listas ---------------------------------------------------------
\setlist{itemsep=2pt, topsep=4pt}

% --- Código (JSON / mongosh) ----------------------------------------
\lstdefinestyle{codigo}{
  basicstyle=\ttfamily\footnotesize,
  breaklines=true,
  columns=fullflexible,
  keepspaces=true,
  frame=single,
  framerule=0.3pt,
  rulecolor=\color{gray!50},
  backgroundcolor=\color{gray!6},
  xleftmargin=4pt, xrightmargin=4pt,
  showstringspaces=false,
  literate={á}{{\'a}}1 {é}{{\'e}}1 {í}{{\'i}}1 {ó}{{\'o}}1 {ú}{{\'u}}1
           {Á}{{\'A}}1 {É}{{\'E}}1 {Í}{{\'I}}1 {Ó}{{\'O}}1 {Ú}{{\'U}}1
           {ñ}{{\~n}}1 {Ñ}{{\~N}}1
}
\lstset{style=codigo}

% --- Atajos ---------------------------------------------------------
\newcommand{\campo}[1]{\texttt{#1}}

\begin{document}

% ------------------------------- Portada -----------------------------
\begin{titlepage}
  \centering
  {\large Pontificia Universidad Javeriana\par}
  {\large Departamento de Ingeniería de Sistemas\par}
  \vspace{0.4cm}
  {\large Administración de Bases de Datos (DBA) -- 2026-30\par}
  \vspace{3cm}
  {\Huge\bfseries Proyecto final\par}
  \vspace{0.4cm}
  {\LARGE Entrega 1: Diseño del esquema NoSQL\\ y plan de implementación\par}
  \vspace{0.8cm}
  {\large Potencial solar en techos de los municipios PDET\par}
  \vfill
  {\large
    Maria Camila Beltran \\
    Jorge Fortich \\
    Jose Eduardo Valeriano \\
    Samuel Osorio Rodríguez\par}
  \vspace{1cm}
  {\large Profesor: Andrés Oswaldo Calderón Romero, Ph.D.\par}
  \vspace{0.4cm}
  {\large 2 de noviembre de 2026\par}
\end{titlepage}

\tableofcontents
\newpage

% ------------------------------- Secciones ---------------------------
% =====================================================================
% 1. Contexto y objetivo
% =====================================================================
\section{Contexto y objetivo}
\label{sec:contexto}

% =====================================================================
% 2. Decisiones tecnológicas
% =====================================================================
\section{Decisiones tecnológicas}
\label{sec:decisiones}

% =====================================================================
% 3. Plan de implementación
% =====================================================================
\section{Plan de implementación}
\label{sec:plan}

\begin{enumerate}[leftmargin=*]
  \item \textbf{Infraestructura (semana 1)}

  \item \textbf{Límites PDET (semana 2)}

  \item \textbf{Huellas de edificios (semana 3)}

  \item \textbf{Análisis espacial (semana 4)}

  \item \textbf{Resultados e informe (semana 4)}

\end{enumerate}

% =====================================================================
% 4. Diseño del esquema
% =====================================================================
\section{Diseño del esquema}
\label{sec:esquema}

% =====================================================================
% 5. Diccionario de datos
% =====================================================================
\section{Diccionario de datos}
\label{sec:diccionario}

\subsection{\texorpdfstring{\campo{municipios}}{municipios}}

\subsection{\texorpdfstring{\campo{edificios\_google} y \campo{edificios\_microsoft}}{edificios\_google y edificios\_microsoft}}

\subsection{\texorpdfstring{\campo{fuentes}}{fuentes}}

\subsection{\texorpdfstring{\campo{resultados}}{resultados}}

\subsection{\texorpdfstring{\campo{bitacora\_cargas}}{bitacora\_cargas}}

\subsection{Índices}

% ------------------------------- Referencias -------------------------
% =====================================================================
% Referencias
% =====================================================================
\begin{thebibliography}{9}

\bibitem{pdet}
Agencia de Renovación del Territorio, \emph{Municipios PDET}
(\texttt{MunicipiosPDET.xlsx}, carpeta de recursos del curso).
\url{https://www.renovacionterritorio.gov.co/central-pdet/que-son-los-pdet}

\bibitem{mgn}
DANE, \emph{Marco Geoestadístico Nacional -- Descarga de datos
geoestadísticos}, 2025.
\url{https://geoportal.dane.gov.co/servicios/descarga-y-metadatos/datos-geoestadisticos/?cod=111}

\bibitem{google}
Google Research, \emph{Open Buildings V3 Polygons}, catálogo de Earth Engine.
\url{https://developers.google.com/earth-engine/datasets/catalog/GOOGLE_Research_open-buildings_v3_polygons}

\bibitem{microsoft}
Microsoft, \emph{ms-buildings: Global ML Building Footprints}, Planetary
Computer.
\url{https://planetarycomputer.microsoft.com/dataset/ms-buildings}

\bibitem{mongodb}
MongoDB, Inc., \emph{Query a 2dsphere Index}, documentación de MongoDB.
\url{https://www.mongodb.com/docs/current/tutorial/query-a-2dsphere-index/}

\end{thebibliography}

\end{document}